The paired-condition design of \textsc{AnchorBench} (identical
prompts differing only in the anchor-bearing component) is
well-suited for mechanistic interpretability analyses such as
logit-lens probing and activation patching.
While mechanistic analysis is not the primary contribution of
this benchmark paper, we briefly report preliminary findings
from a v1-benchmark study to illustrate the benchmark's
utility for this purpose.

\paragraph{Setup.}
Using 200 stratified items from the v1 benchmark (which used
the plausible relevance level only), we performed logit-lens
and residual-stream patching analyses on three open-weight
models: Llama-3.2-1B, Llama-3.2-3B, and Qwen3-4B-FP8.

\paragraph{Key observations.}
(1)~The anchor--control JS divergence remains near zero in
early transformer layers and rises sharply in the final third
of the network, indicating that anchoring distortion is
introduced or amplified in late layers.
(2)~Attention-head attribution reveals architecture-dependent
patterns: Qwen3-4B shows concentrated anchoring effects in a
small number of heads (head~0 in layers 20--27, effect sizes
0.40--0.46), while the Llama models show distributed
attribution with smaller individual effects (0.06--0.07).
(3)~These findings are preliminary and based on models at
1B--4B scale; generalization to larger models and other
architectures is an open question.

\paragraph{Benchmark utility.}
These results demonstrate that \textsc{AnchorBench}'s
paired-condition design enables clean causal attribution of
anchoring effects---a capability that, combined with the
relevance axis, could support future investigations into
whether models process irrelevant and plausible anchors through
different internal pathways.
